\documentclass[11pt,a4paper]{article}
% Préambule commun aux documents de PythonIPM/documentation.
% Compilation : tectonic <fichier>.tex
% tectonic compile avec XeTeX, qui lit l'UTF-8 nativement : surtout PAS inputenc ni fontenc,
% faits pour pdfLaTeX. Ce sont eux qui décodaient mal les caractères multi-octets — « cœur »
% ressortait en « cur », et les tirets cadratins en caractères parasites.
\usepackage{lmodern}
\usepackage[french]{babel}
\usepackage{microtype}
\usepackage[a4paper,margin=2.4cm]{geometry}
\usepackage{amsmath,amssymb}
\usepackage{booktabs}
\usepackage{longtable}
\usepackage{array}
\usepackage{ragged2e}
\usepackage{xcolor}
\usepackage{tcolorbox}
\usepackage{enumitem}
\usepackage{fancyhdr}
\usepackage[hidelinks]{hyperref}

\definecolor{encadre}{HTML}{F4F6F8}
\definecolor{bordure}{HTML}{4A6FA5}
\definecolor{alerte}{HTML}{B03A2E}

% Encadré « à retenir » : la lecture non technique du passage qui précède
\newtcolorbox{retenir}[1][]{
  colback=encadre, colframe=bordure, boxrule=0.6pt, arc=2pt,
  left=8pt, right=8pt, top=6pt, bottom=6pt,
  fonttitle=\bfseries, title={#1}
}

% Encadré d'avertissement : un point ouvert, un écart assumé
\newtcolorbox{attention}[1][]{
  colback=white, colframe=alerte, boxrule=0.6pt, arc=2pt,
  left=8pt, right=8pt, top=6pt, bottom=6pt,
  % le bandeau de titre a déjà la couleur d'alerte en fond : un titre `coltitle=alerte`
  % écrirait du rouge sur du rouge. On garde le blanc par défaut, comme l'encadré « retenir ».
  fonttitle=\bfseries, title={#1}
}

\newcommand{\col}[1]{\texttt{#1}}
% \ensuremath et non \textsubscript : \ipm est employé aussi bien dans le texte que dans les
% formules, et « M\textsubscript{0} » placé en mode mathématique est une erreur — silencieuse en
% PDF, mais qui fait échouer la conversion vers Word.
\newcommand{\ipm}{\ensuremath{M_{0}}}

\setlist[itemize]{itemsep=2pt, topsep=4pt}
\setlist[description]{style=nextline, leftmargin=1.2cm, itemsep=5pt}

% La source du document apparaît en tête de chaque page. Les documents RGPH la redéfinissent
% par \renewcommand{\source}{...} AVANT \begin{document}.
\newcommand{\source}{EHCVM 2021}

\pagestyle{fancy}
\fancyhf{}
\fancyhead[L]{\small\itshape IPM Côte d'Ivoire --- \source}
\fancyhead[R]{\small\thepage}
\renewcommand{\headrulewidth}{0.3pt}

\renewcommand{\arraystretch}{1.15}


\title{\bfseries L'indice de pauvreté multidimensionnelle de la Côte d'Ivoire\\[4pt]
\large Méthodologie de calcul sur l'EHCVM 2021}
\author{Pipeline \col{PythonIPM}}
\date{Août 2026}

\begin{document}
\maketitle

\begin{abstract}
\noindent Ce document décrit, de bout en bout, la construction de l'indice de pauvreté
multidimensionnelle (IPM) de la Côte d'Ivoire à partir de l'Enquête harmonisée sur les conditions
de vie des ménages (EHCVM) 2021. Il est écrit pour être lu à deux niveaux : le fil du texte
n'exige aucune connaissance préalable, tandis que les formules et les tableaux donnent au lecteur
averti tout ce qu'il faut pour reproduire ou contester le calcul. Chaque choix de définition, chaque
écart à la norme internationale et chaque convention de traitement y sont explicités.
\end{abstract}

\begin{retenir}[Les résultats en cinq lignes]
\begin{itemize}
  \item \textbf{58,0~\%} de la population ivoirienne vit dans un ménage multidimensionnellement
        pauvre, soit 17,3 millions de personnes --- intervalle de confiance à 95~\% :
        [54,9~\% ; 61,0~\%].
  \item Ces personnes cumulent en moyenne \textbf{49,4~\%} des privations pondérées possibles.
  \item L'indice qui combine les deux vaut $\ipm = \mathbf{0{,}2863}$, IC 95~\%
        [0,2701 ; 0,3026].
  \item Le milieu rural rassemble \textbf{71,9~\%} de la pauvreté pour 47,3~\% de la population.
  \item L'écart entre régions va de 0,058 (Abidjan) à 0,481 (Bounkani), soit un rapport de
        un à huit.
  \item Sur le périmètre comparable au niveau international --- les trois dimensions sans
        l'Emploi --- l'indice vaut $\ipm = \mathbf{0{,}3461}$ pour $H = 66{,}4~\%$
        (section~\ref{sec:international}).
\end{itemize}
\end{retenir}

\tableofcontents
\bigskip

\section{Ce que mesure un IPM}

La pauvreté monétaire se mesure avec une seule question : combien un ménage dépense-t-il par
rapport à un seuil ? C'est simple, mais incomplet. Un ménage peut dépasser le seuil monétaire et
n'avoir ni eau potable, ni électricité, ni enfant scolarisé. Inversement, un ménage modeste peut
bénéficier d'une école et d'un dispensaire proches.

L'indice de pauvreté multidimensionnelle répond à cette limite en observant simultanément
plusieurs conditions de vie --- l'éducation, la santé, l'emploi, le logement --- et en comptant
comme pauvre le ménage qui cumule suffisamment de privations.

\begin{retenir}[En une phrase]
L'IPM compte les gens qui cumulent des manques, et mesure à quel point ils en cumulent.
\end{retenir}

La méthode employée est celle d'Alkire et Foster, adoptée par le PNUD et l'OPHI, et reprise dans
le guide « L'indice de pauvreté multidimensionnelle au service des ODD » (chapitre 4). Elle repose
sur une idée simple, dite du \emph{double seuil} :

\begin{enumerate}
  \item un premier seuil, propre à chaque indicateur, dit si le ménage est privé sur ce point
        précis (par exemple : moins de dix années d'études pour le membre le plus instruit) ;
  \item un second seuil, appliqué au total pondéré des privations, dit si le ménage est
        \emph{multidimensionnellement pauvre}.
\end{enumerate}

Le second seuil est ce qui distingue l'IPM d'un simple tableau de bord d'indicateurs : il produit
un chiffre unique, décomposable, dont on peut suivre l'évolution et attribuer les variations à
telle région ou tel indicateur.

\section{Les données}

L'EHCVM 2021 de la Côte d'Ivoire couvre \textbf{12\,965 ménages}, soit
\textbf{29\,809\,416 personnes} une fois les pondérations appliquées. Elle a été collectée en deux
vagues, la première fin 2021, la seconde début 2022. Quatre bases sur les vingt-et-une du dossier
d'enquête portent des indicateurs utiles à l'IPM.

\begin{center}
\begin{tabular}{@{}llr>{\RaggedRight}p{6.2cm}@{}}
\toprule
\textbf{Base} & \textbf{Une ligne =} & \textbf{Lignes} & \textbf{Ce qu'on y prend} \\
\midrule
\col{Base\_Individus} & un membre du ménage & 64\,491 & éducation, emploi, état civil, santé \\
\col{Base\_Menage} & un ménage & 12\,965 & conditions de vie, pondération, désagrégation \\
\col{Base\_avoirs\_du\_menage} & un bien proposé & 583\,425 & biens d'équipement \\
\col{Base\_securite\_alimentaire} & un ménage & 13\,693 & insécurité alimentaire (FIES) \\
\bottomrule
\end{tabular}
\end{center}

Les quatre bases se raccordent par la clé \col{grappe} + \col{menage} + \col{vague} et couvrent
exactement les mêmes ménages : la fusion se fait sans perte. La \col{grappe} est la zone de
dénombrement tirée au premier degré du sondage (un quartier, un village) ; \col{menage} est le
numéro d'ordre du ménage dans sa grappe ; \col{vague} distingue les deux passages.

\subsection{La pondération, et pourquoi elle est indispensable}

L'échantillon est volontairement déséquilibré : on interroge proportionnellement plus de ménages
dans les régions peu peuplées, afin d'y disposer d'assez d'observations. La variable
\col{ponderation\_menage} (\col{hhweight} dans le fichier d'origine) corrige ce déséquilibre : elle
indique combien de ménages réels chaque ménage enquêté représente.

L'ordre de grandeur de la correction n'a rien d'anecdotique : le milieu urbain représente
\textbf{40,4~\%} des lignes du fichier mais \textbf{52,7~\%} de la population une fois pondéré.
Sans pondération, tous les indices seraient biaisés vers le rural, donc l'IPM surestimé.

\begin{retenir}[Deux « pondérations » à ne pas confondre]
\begin{itemize}
  \item la \textbf{pondération d'enquête} (\col{ponderation\_menage}) sert à passer de
        l'échantillon à la population ivoirienne ;
  \item les \textbf{poids $w_j$} de la méthode Alkire-Foster disent l'importance relative de
        chaque indicateur dans le score d'un ménage.
\end{itemize}
Elles portent le même nom courant et n'ont rien à voir.
\end{retenir}

\subsection{Unité de construction et unité de comptage}

Le score de privation se calcule \textbf{une fois par ménage} : c'est le ménage qui a un toit, une
source d'eau, des toilettes. Mais l'indice se lit \textbf{en personnes} : un ménage de huit
personnes qui cumule les privations pèse davantage qu'un ménage d'une personne.

Chaque ménage $i$ compte donc pour
\[
n_i \;=\; \col{taille\_menage}_i \times \col{ponderation\_menage}_i .
\]

\section{Les seize indicateurs}

Quatre dimensions, équipondérées à 0,25 chacune. Le poids d'une dimension se partage à parts
égales entre ses indicateurs : deux indicateurs portent 0,125 chacun, quatre indicateurs
0,0625 chacun.

\begin{center}
\begin{tabular}{@{}lrrr@{}}
\toprule
\textbf{Dimension} & \textbf{Indicateurs} & \textbf{Poids $w_j$} & \textbf{Total} \\
\midrule
Éducation & 4 & 0,0625 & 0,25 \\
Santé & 3 & 0,0833 & 0,25 \\
Emploi & 2 & 0,1250 & 0,25 \\
Conditions de vie & 7 & 0,0357 & 0,25 \\
\midrule
\textbf{Ensemble} & \textbf{16} & & \textbf{1,00} \\
\bottomrule
\end{tabular}
\end{center}

Chaque indicateur suit \emph{soit} la proposition nationale, \emph{soit} l'application du PNUD.
Ce choix a été arrêté indicateur par indicateur ; il figure dans le code et dans le fichier
\col{vecteur\_z.csv} produit à chaque exécution.

\begin{center}
\small
\begin{longtable}{@{}>{\RaggedRight}p{2.1cm}>{\RaggedRight}p{2.5cm}l>{\RaggedRight}p{5.4cm}r@{}}
\toprule
\textbf{Dimension} & \textbf{Indicateur} & \textbf{Source} & \textbf{Un ménage est privé si\ldots} &
\textbf{Privés} \\
\midrule
\endfirsthead
\toprule
\textbf{Dimension} & \textbf{Indicateur} & \textbf{Source} & \textbf{Un ménage est privé si\ldots} &
\textbf{Privés} \\
\midrule
\endhead
Éducation & Fréquentation scolaire & nationale & un enfant de 6--16 ans n'est pas scolarisé
& 27,9~\% \\
Éducation & Année de scolarité & nationale & aucun membre de 17--95 ans n'a atteint 10 années
d'études & 72,5~\% \\
Éducation & Alphabétisation & nationale & un membre de 17--49 ans ne sait ni lire ni écrire le
français & 61,7~\% \\
Éducation & État civil & nationale & un enfant de 5--15 ans n'a pas d'acte de naissance
& 23,5~\% \\
\midrule
Santé & Assurance maladie & nationale & aucun membre n'est couvert par une assurance maladie
& 93,4~\% \\
Santé & Insécurité alimentaire & FAO & le score FIES atteint 4 sur 8 (insécurité modérée ou
sévère) & 42,7~\% \\
Santé & Renoncement aux soins (coût ou indisponibilité) & nationale & un membre malade dans les
30 derniers jours n'a pas consulté, pour raison de coût ou faute de service disponible
& 9,3~\% \\
\midrule
Emploi & Chômage & nationale & un membre de 17--40 ans est au chômage au sens du BIT & 3,7~\% \\
Emploi & Emploi agricole de subsistance & nationale & le chef de ménage occupé ne travaille que
sur son propre champ : ni salariat, ni apprentissage, ni commerce & 48,1~\% \\
\midrule
Cadre de vie & Électricité & nationale & l'éclairage n'est ni le réseau, ni un groupe
électrogène, ni le solaire & 13,3~\% \\
Cadre de vie & Logement & PNUD & le sol, le toit ou les murs sont en matériaux naturels ou
rudimentaires & 22,3~\% \\
Cadre de vie & Eau potable & PNUD & la source n'est pas améliorée, ou est à plus de 15 minutes
à l'aller & 24,6~\% \\
Cadre de vie & Énergie de cuisson & nationale & le combustible principal n'est ni le gaz ni
l'électricité & 76,8~\% \\
Cadre de vie & Toilettes & PNUD & les sanitaires ne sont pas améliorés, ou sont partagés
& 77,3~\% \\
Cadre de vie & Biens d'équipement & PNUD & le ménage possède au plus un bien sur sept et pas de
voiture & 29,2~\% \\
Cadre de vie & Promiscuité & nationale & le logement abrite plus de 3 personnes par pièce
(seuil ONU-Habitat) & 8,4~\% \\
\bottomrule
\end{longtable}
\end{center}

\begin{retenir}[« Non concerné » n'est pas « privé »]
Un ménage sans enfant de 6 à 16 ans n'est pas privé de fréquentation scolaire : la question ne le
concerne pas. C'est la convention retenue, appliquée uniformément. Le pipeline conserve, à côté de
chaque effectif privé, l'effectif \emph{concerné}, de sorte que cette convention peut être
renversée sans revenir aux données individuelles.
\end{retenir}

\section{Le calcul, étape par étape}

\subsection{La matrice situationnelle et le vecteur de seuils}

On part d'un tableau $X$ à 12\,965 lignes (les ménages) et 16 colonnes (les indicateurs), qui
décrit la \emph{situation} de chaque ménage : combien d'enfants non scolarisés, quel type de
sanitaires, combien d'années d'études pour le plus instruit.

Le vecteur $z = (z_1, \dots, z_{16})$ donne le seuil de chaque indicateur. Il ne prend pas
toujours la forme d'un nombre : selon l'indicateur, le seuil est une valeur numérique (10 années
d'études), un effectif (« au moins une personne »), ou un ensemble de modalités jugées adéquates
(les codes des sources d'eau améliorées au sens des ODD). Ces trois formes coexistent dans
\col{vecteur\_z.csv}.

\subsection{La matrice de privation $g^0$}

Appliquer $z$ à $X$ donne une matrice de $0$ et de $1$ :
\[
g^0_{ij} =
\begin{cases}
1 & \text{si le ménage } i \text{ est privé sur l'indicateur } j,\\
0 & \text{sinon.}
\end{cases}
\]

\subsection{La pondération et le score $c_i$}

Chaque colonne est multipliée par son poids, puis la ligne est sommée :
\[
c_i \;=\; \sum_{j=1}^{16} w_j \, g^0_{ij}, \qquad 0 \le c_i \le 1 .
\]

Le score $c_i$ est la \emph{part pondérée de privations} du ménage $i$. Un ménage privé partout a
$c_i = 1$ ; un ménage sans aucune privation a $c_i = 0$. Sur l'EHCVM 2021, le score moyen est de
\textbf{0,3918} et le maximum observé de 0,8929 : aucun ménage n'est privé sur les seize
indicateurs à la fois.

\subsection{La censure au seuil \texorpdfstring{$k$}{k}}

Un ménage est déclaré multidimensionnellement pauvre si son score atteint $k = 1/3$, c'est-à-dire
s'il cumule l'équivalent d'une dimension et un tiers de privations. On définit alors le
\emph{score censuré} :
\[
c_i(k) =
\begin{cases}
c_i & \text{si } c_i \ge k,\\
0 & \text{sinon.}
\end{cases}
\]

La censure est le cœur de la méthode. Les privations des ménages \emph{non} pauvres ne
disparaissent pas des données --- elles restent lisibles dans la matrice pondérée --- mais elles
ne comptent pas dans l'indice. Sur l'EHCVM 2021, la censure efface \textbf{18,3~\%} des privations
pondérées, celles de 4\,269 ménages.

\begin{retenir}[Pourquoi censurer ?]
Sans censure, une amélioration chez un ménage déjà très peu privé ferait baisser l'indice autant
qu'une amélioration chez un ménage très pauvre. La censure garantit que l'IPM ne parle que des
pauvres, et donc qu'une politique qui les cible fait bouger l'indice.
\end{retenir}

Deux seuils complémentaires sont publiés à côté : la \textbf{vulnérabilité}
($0{,}2 \le c_i < 1/3$) et la \textbf{pauvreté sévère} ($c_i \ge 0{,}5$).

\subsection{Les trois indices}

\begin{align}
H  &= \frac{\sum_i n_i \, \mathbf{1}(c_i \ge k)}{\sum_i n_i}
   && \text{\emph{incidence} : la part de la population pauvre} \\[4pt]
A  &= \frac{\sum_i n_i \, c_i(k)}{\sum_i n_i \, \mathbf{1}(c_i \ge k)}
   && \text{\emph{intensité} : les privations moyennes des pauvres} \\[4pt]
\ipm &= H \times A \;=\; \frac{\sum_i n_i \, c_i(k)}{\sum_i n_i}
   && \text{l'IPM}
\end{align}

$H$ répond à « combien de pauvres ? », $A$ à « à quel point sont-ils pauvres ? », et \ipm{} combine
les deux. Un pays peut faire baisser $H$ en sortant de la pauvreté les ménages qui en étaient les
plus proches, sans rien changer pour les plus démunis : \ipm{} le détecte, $H$ seul non.

\subsection{Les contributions}

L'IPM est \emph{décomposable} par indicateur. La contribution de l'indicateur $j$ est
\[
C_j = \frac{w_j \, CH_j}{\ipm}, \qquad \sum_j C_j = 1,
\]
où $CH_j$ est le \emph{taux de privation censuré} : la proportion de la population qui est à la
fois pauvre et privée sur l'indicateur $j$. Il est également décomposable par sous-population :
la somme des \ipm{} de groupe, pondérée par leur part de population, redonne exactement l'\ipm{}
national. Cette propriété est vérifiée automatiquement à chaque exécution du pipeline.

\section{Résultats}

\begin{center}
\begin{tabular}{@{}lrll@{}}
\toprule
\textbf{Indice} & \textbf{Valeur} & \textbf{IC 95~\%} & \textbf{Lecture} \\
\midrule
$H$ --- incidence & 0,5796 & [0,5490 ; 0,6102] & 58,0~\% de la population est pauvre, soit
17,3 millions \\
$A$ --- intensité & 0,4940 & [0,4894 ; 0,4986] & les pauvres cumulent 49,4~\% des privations
pondérées \\
\ipm{} & 0,2863 & [0,2700 ; 0,3026] & $= H \times A$ \\
Vulnérables & 22,8~\% & & score entre 0,2 et 1/3 \\
Pauvreté sévère & 26,1~\% & & score d'au moins 0,5 \\
\bottomrule
\end{tabular}
\end{center}

\subsection{Où se concentre la pauvreté}

\begin{center}
\begin{tabular}{@{}lrrrr@{}}
\toprule
& \textbf{Part de la population} & \textbf{$H$} & \textbf{$A$} & \textbf{\ipm{}} \\
\midrule
Urbain & 52,7~\% & 33,6~\% & 0,4533 & 0,1525 \\
Rural  & 47,3~\% & 85,1~\% & 0,5120 & 0,4354 \\
\midrule
Chef de ménage homme & 83,6~\% & 60,7~\% & 0,4970 & 0,3014 \\
Chef de ménage femme & 16,4~\% & 44,2~\% & 0,4733 & 0,2092 \\
\bottomrule
\end{tabular}
\end{center}

Le rural concentre \textbf{71,9~\%} de la pauvreté multidimensionnelle avec 47,3~\% de la
population ; son \ipm{} est près de trois fois celui de l'urbain. Entre régions, l'écart va de
0,058 (Abidjan) à 0,481 (Bounkani). Trois régions --- Worodougou, Béré, Folon --- dépassent
90~\% d'incidence.

Le résultat apparemment contre-intuitif --- les ménages dirigés par une femme sont \emph{moins}
pauvres --- est classique en IPM : ces ménages sont en moyenne plus petits et plus urbains. Il ne
dit rien du niveau de vie individuel des femmes, que l'IPM ménage ne mesure pas.

\subsection{Ce que pèse réellement chaque dimension}

\begin{center}
\begin{tabular}{@{}lrr@{}}
\toprule
\textbf{Dimension} & \textbf{Poids nominal} & \textbf{Contribution réelle à \ipm{}} \\
\midrule
Éducation & 25,0~\% & 32,3~\% \\
Santé & 25,0~\% & 27,9~\% \\
Conditions de vie & 25,0~\% & 21,3~\% \\
Emploi & 25,0~\% & 18,5~\% \\
\bottomrule
\end{tabular}
\end{center}

\begin{attention}[Un résultat qui mérite l'attention du lecteur]
Les quatre dimensions sont équipondérées \emph{en droit} ; elles pèsent en fait de 18,5~\% à
32,3~\%. Le poids $w_j$ fixe le maximum qu'une dimension \emph{peut} apporter ; ce qu'elle
apporte \emph{vraiment} dépend de la fréquence de la privation. L'écart s'est nettement resserré
depuis la version à treize indicateurs, où l'Emploi tombait à 4,6~\% et la Santé montait à
40,5~\% : ces deux dimensions n'avaient alors qu'un ou deux indicateurs, et se trouvaient donc à
la merci de la fréquence de chacun. Le déséquilibre s'est déplacé \emph{à l'intérieur} des
dimensions : en Emploi, l'agriculture de subsistance apporte 16,9~\% de l'\ipm{} et le chômage
1,6~\%, à poids égal (0,125) --- le premier prive 48,1~\% des ménages, le second 3,7~\%.
\end{attention}

\section{La précision des estimations}

Tous les indices publiés sont des estimations tirées d'un échantillon : ils sont accompagnés
d'un \textbf{intervalle de confiance à 95~\%}, calculé pour $H$, $A$ et \ipm{}, à tous les
niveaux de désagrégation.

\subsection{Le plan de sondage, pas la loi binomiale}

$H$, $A$ et \ipm{} sont tous des ratios de la même forme
\[
R = \frac{\sum_i n_i \, y_i}{\sum_i n_i \, x_i}
\]
--- pour $H$, $y_i = \mathbf{1}(c_i \ge k)$ et $x_i = 1$ ; pour \ipm{}, $y_i = c_i(k)$ et
$x_i = 1$ ; pour $A$, $y_i = c_i(k)$ et $x_i = \mathbf{1}(c_i \ge k)$. Une seule formule de
variance les couvre donc tous les trois.

L'EHCVM est un sondage \textbf{à deux degrés} : on tire d'abord des grappes (zones de
dénombrement), puis des ménages dans chaque grappe. Deux ménages d'une même grappe se
ressemblent --- même village, même école, même point d'eau --- et n'apportent donc pas deux
informations indépendantes. La variance est calculée en conséquence, par agrégation des
résidus linéarisés \textbf{au niveau de la grappe} :
\[
\widehat{\mathrm{V}}(R) = \frac{m}{m-1} \sum_{g=1}^{m} u_g^2,
\qquad
u_g = \sum_{i \in g} \frac{n_i \,(y_i - R\,x_i)}{\sum_i n_i x_i},
\]
où $m$ est le nombre de grappes du domaine. L'intervalle est
$R \pm t_{0{,}975;\,m-1}\sqrt{\widehat{\mathrm{V}}(R)}$, borné à $[0, 1]$. Traiter les ménages
comme indépendants diviserait l'erreur-type par deux ou plus : les intervalles publiés seraient
alors faux, et faussement rassurants.

\begin{retenir}[Ce que l'intervalle recouvre --- et ce qu'il ne recouvre pas]
L'intervalle décrit la seule \textbf{incertitude d'échantillonnage}. Il ne dit rien des erreurs
de déclaration, des non-réponses, ni des choix de définition documentés dans ce rapport --- qui
pèsent souvent bien plus lourd. Un \ipm{} de 0,2863 « à $\pm$ 0,016 près » reste un \ipm{} dont
la valeur dépend d'abord des seize indicateurs retenus.
\end{retenir}

\subsection{Jusqu'où descendre dans le territoire}

Les indices sont désagrégés à quatre niveaux : milieu, région (33), \textbf{département} (108) et
\textbf{sous-préfecture ou commune} (442). Plus le domaine est fin, moins il contient de grappes,
et plus l'intervalle s'élargit.

\begin{center}
\begin{tabular}{@{}lrrrr@{}}
\toprule
\textbf{Niveau} & \textbf{Domaines} & \textbf{Grappes (médiane)} & \textbf{Largeur de l'IC} &
\textbf{CV $>$ 20~\%} \\
\midrule
Milieu & 2 & 540 & 0,027 & 0 \\
Région & 33 & 30 & 0,113 & 0 \\
Département & 108 & 8 & 0,213 & 15 \\
Sous-préfecture / commune & 442 & 2 & 0,485 & 92 \\
\bottomrule
\end{tabular}
\end{center}

La colonne \col{M0\_cv} du classeur donne le coefficient de variation de \ipm{}
(erreur-type / estimation). L'usage statistique retient 20~\% comme limite de publication : au
delà, l'estimation est indicative et ne doit pas être commentée seule. Au niveau régional, aucune
estimation n'est concernée ; au niveau des sous-préfectures, 92 sur 442 le sont, et
\textbf{211 domaines ne comptent qu'une seule grappe}, ce qui rend la variance non estimable ---
leurs colonnes d'intervalle sont alors vides, jamais remplies d'une valeur inventée.

\begin{attention}[Le département est le plancher raisonnable]
L'EHCVM 2021 a été dimensionnée pour être représentative au niveau régional. Les résultats
départementaux sont exploitables avec prudence ; les résultats par sous-préfecture sont fournis
pour le ciblage local et l'exploration, pas pour la publication ni pour classer des territoires
entre eux.
\end{attention}

\section{Écarts assumés et limites}

L'EHCVM n'a pas été conçue pour calculer un IPM. Plusieurs indicateurs de la norme internationale
ont dû être adaptés. Ces écarts sont documentés ici plutôt que dissimulés : ils conditionnent
l'interprétation des résultats et la comparabilité internationale.

\begin{description}
  \item[Mortalité juvénile --- indicateur abandonné.] L'enquête n'a ni module décès ni historique
  des naissances. La base des chocs enregistre bien un décès dans le ménage (1\,328 ménages sur
  trois ans) mais sans l'âge du défunt : impossible d'isoler les moins de 18 ans. La dimension
  Santé repose donc sur l'assurance maladie, l'insécurité alimentaire et le renoncement aux soins. Ce dernier indicateur
  remplace la nutrition de l'IPM mondial, que l'EHCVM ne permet pas non plus de mesurer faute de
  module anthropométrique. Cette substitution fait l'objet d'un document séparé.

  \item[Années d'études --- reconstruites.] L'EHCVM ne demande pas le nombre d'années d'études.
  Il est reconstitué comme « années accomplies avant le niveau + classe atteinte », en trois
  branches : niveau achevé, niveau en cours (moins l'année non terminée), et zéro pour ceux qui
  n'ont jamais fréquenté l'école. Restent 8,5~\% d'individus indéterminés, qui ne concernent que
  six ménages une fois l'agrégation faite par le maximum.

  \item[Charrette --- absente.] Les 45 biens de la section 12 de l'EHCVM ne comprennent ni
  charrette ni camion. Le décompte des biens d'équipement porte donc sur \emph{sept} biens au lieu
  de huit, ce qui rend le ménage marginalement moins souvent privé.

  \item[Chômage --- reconstitué.] Aucune variable de statut d'activité au sens du BIT n'existe
  telle quelle. Les trois critères sont recomposés à partir de six questions : sans emploi, en
  recherche, et disponible.

  \item[Fréquentation scolaire --- deux années scolaires.] La question portant sur 2021/22 n'est
  posée qu'en seconde vague. Un enfant est réputé scolarisé s'il l'était sur l'année la plus
  récente disponible pour sa vague.

  \item[Renoncement aux soins --- conditionné à la maladie, et élargi à l'offre.] L'indicateur ne
  peut être privé que dans un ménage où quelqu'un a été malade dans les 30 derniers jours :
  67,2~\% des ménages sont concernés, et 9,3~\% sont privés. Un ménage où personne n'a été malade
  n'est pas privé, comme partout ailleurs dans ce pipeline. Quatre modalités de la question 3.06
  rendent privé, réunies par le fait d'être \emph{subies} : le coût --- 2 (« trop cher ») et 8
  (« manque d'argent ») --- et l'indisponibilité de l'offre --- 11 (« service spécialisé non
  disponible ») et 12 (« absence de personnel »). Les motifs d'arbitrage, au premier rang
  desquels l'automédication, ne rendent pas privé. L'élargissement à l'offre ne concerne que 14
  personnes et 10 ménages : il pèse 0,00008 sur l'\ipm{} et relève de la définition, non du
  résultat.

  \item[Emploi de subsistance --- codage 1/2, pas 0/1.] Les questions 4.06 à 4.09 de l'EHCVM sont
  codées 1 = oui, 2 = non ; elles ne prennent jamais la valeur 0. Le critère « sans emploi
  salarié, ni apprentissage, ni commerce » est donc appliqué comme « pas de \emph{oui} », ce qui
  englobe les non-réponses de ces branches. L'indicateur porte sur le seul \textbf{chef de
  ménage} : c'est lui qui détermine, dans la structure de l'EHCVM, l'activité principale du
  ménage.

  \item[Promiscuité --- seuil ONU-Habitat.] Le seuil retenu est de \emph{plus de} 3 personnes par
  pièce (question 11.02), recommandé par ONU-Habitat pour l'Afrique subsaharienne et utilisé par
  les IPM nationaux éthiopien et rwandais : 8,4~\% des ménages sont privés. Le seuil de l'OMS,
  plus exigeant, est de 2 personnes par pièce et porterait la privation à 29,7~\%. La bascule
  tient à une constante du code.

  \item[Électricité --- source d'éclairage, pas raccordement.] La variable « connecté au réseau »
  ignore le groupe électrogène et le solaire : elle classerait 2\,597 ménages de plus comme privés
  (33,3~\% contre 13,3~\%). L'indicateur repose sur la source d'éclairage effective.
\end{description}

\begin{attention}[Deux points restés ouverts]
\textbf{Alphabétisation.} L'énoncé national dit qu'un ménage est privé si \emph{un} membre de
17--49 ans ne sait ni lire ni écrire (61,7~\% des ménages). Le recensement de 2021 codait
l'inverse : privé si \emph{aucun} membre n'est alphabétisé (35,0~\%). Le code suit l'énoncé
littéral ; la bascule tient à une constante.

\medskip
\textbf{Seuil de l'eau.} La norme du PNUD dit « 30 minutes ou plus » aller-retour, ce qui
correspond à 15 minutes ou plus à l'aller, la question de l'EHCVM mesurant le trajet aller seul.
Le code applique un seuil strict (\emph{plus de} 15 minutes). L'écart porte sur 402 ménages qui
déclarent exactement 15 minutes --- un effet d'arrondi aux multiples de cinq --- et ferait passer
la privation en eau de 24,6~\% à environ 27,7~\%.
\end{attention}

\section{La variante internationale : trois dimensions}
\label{sec:international}

Les adaptations qui précèdent, et surtout le choix des dimensions, interdisent une comparaison
directe avec l'IPM mondial publié par l'OPHI, qui repose sur trois dimensions et dix indicateurs.
La dimension \textbf{Emploi} est en particulier une spécificité nationale : elle n'existe ni dans
l'IPM mondial ni dans le tableau de référence ODD.

Le pipeline publie donc \textbf{deux indices}, calculés par le même code sur les mêmes données,
qui ne diffèrent que par le périmètre des dimensions retenues :

\begin{center}
\begin{tabular}{@{}lll@{}}
\toprule
& \textbf{IPM national} & \textbf{IPM international} \\
\midrule
Dimensions & 4 & 3 (sans l'Emploi) \\
Indicateurs & 16 & 14 \\
Poids d'une dimension & 0,25 & 1/3 \\
Table publiée & \col{indices\_ipm\_ci} & \col{indices\_ipm\_international} \\
\bottomrule
\end{tabular}
\end{center}

Retirer l'Emploi ne consiste pas à mettre ses deux indicateurs à zéro : les trois dimensions
restantes sont \textbf{réequipondérées} à 1/3 chacune. Chaque indicateur d'Éducation passe donc
de 0,0625 à 1/12, chaque indicateur de Santé de 0,0833 à 1/9, et chaque indicateur de conditions
de vie de 0,0357 à 1/21. Le seuil $k = 1/3$, la censure et les formules de $H$, $A$ et \ipm{}
sont inchangés.

\begin{center}
\begin{tabular}{@{}lrrr@{}}
\toprule
& \textbf{National (16)} & \textbf{International (14)} & \textbf{Écart} \\
\midrule
$H$ --- incidence & 0,5796 & 0,6639 & $+8{,}4$ pts \\
$A$ --- intensité & 0,4940 & 0,5214 & $+2{,}7$ pts \\
\ipm{} & 0,2863 & 0,3461 & $+0{,}060$ \\
IC 95~\% de \ipm{} & [0,2700 ; 0,3026] & [0,3296 ; 0,3624] & \\
Vulnérables & 22,8~\% & 18,0~\% & $-4{,}8$ pts \\
Pauvreté sévère & 26,1~\% & 34,9~\% & $+8{,}8$ pts \\
Population pauvre & 17,3 millions & 19,8 millions & $+2{,}5$ M \\
\bottomrule
\end{tabular}
\end{center}

\begin{attention}[Retirer une dimension fait \emph{monter} l'indice]
Le résultat surprend, et il est pourtant mécanique. Les deux indicateurs d'Emploi sont peu
fréquents rapportés à leur poids : ensemble ils portent 0,25 du score mais n'apportent que
18,5~\% de l'\ipm{} national. Les retirer redistribue ce quart aux trois dimensions restantes,
où les privations sont bien plus répandues --- l'assurance maladie prive 93,4~\% des ménages,
l'année de scolarité 72,5~\%. Le score de chaque ménage augmente, davantage de ménages franchissent
$k = 1/3$, et $H$ comme $A$ montent. \textbf{Un IPM plus élevé ne signifie donc pas ici une
pauvreté plus grande : c'est le même pays mesuré avec une autre règle.} Les deux chiffres ne se
comparent qu'à eux-mêmes dans le temps, jamais l'un à l'autre.
\end{attention}

Le déséquilibre entre poids nominal et poids réel se creuse dans la variante internationale : à
1/3 nominal, l'Éducation apporte 39,2~\% de l'\ipm{}, la Santé 35,0~\% et les conditions de vie
25,8~\%. L'assurance maladie devient à elle seule le premier contributeur (21,1~\%), là où elle
était au coude à coude avec l'emploi de subsistance dans la variante nationale.

\begin{retenir}[Lequel publier ?]
L'\textbf{IPM national à 16 indicateurs} reste l'indice de référence pour la Côte d'Ivoire : il
intègre l'Emploi, dimension retenue par la proposition nationale. L'\textbf{IPM international}
sert aux comparaisons régionales et au suivi des ODD. Les deux tables ont exactement la même
structure de colonnes et les mêmes désagrégations, ce qui rend leur mise en regard immédiate.
Même ainsi, la comparaison avec l'IPM mondial de l'OPHI reste indicative : les indicateurs de
santé de l'EHCVM ne sont pas ceux de l'OPHI (voir la \emph{Note sur la dimension Santé}).
\end{retenir}

\section{Le traitement des valeurs manquantes}

Le pipeline distingue deux moments, et c'est délibéré.

\textbf{En amont}, la table de préconstruction conserve les valeurs manquantes telles qu'elles
sont dans l'enquête. Quatre colonnes en portent : le temps de trajet vers la source d'eau
(6\,423), le partage des sanitaires (3\,194), le score FIES (195) et les années d'études (18).

\textbf{Au moment d'appliquer le seuil}, une situation restée indéterminée devient « non privé ».
La matrice de privation ne contient donc aucune valeur manquante, ce que deux contrôles
automatiques vérifient à chaque exécution.

Sur les 9\,830 valeurs manquantes, \textbf{9\,617 sont des filtres du questionnaire} où « non
privé » se déduit du filtre lui-même : si l'on n'a pas demandé le temps de trajet, c'est que l'eau
est sur place. Seules \textbf{195} --- les scores FIES incomplets pour cause de non-réponse ---
relèvent d'un véritable choix : celui de ne pas pénaliser un refus de répondre.

\begin{retenir}[Le garde-fou]
Pour chaque indicateur mesuré sur les personnes, la préconstruction conserve deux colonnes :
l'effectif \emph{concerné} et l'effectif \emph{défavorable}. Un zéro dans la seconde ne dit pas
s'il signifie « personne n'est privé » ou « personne n'est concerné » ; la première le dit. Toute
convention peut ainsi être révisée sans repartir des données individuelles.
\end{retenir}

\section{Le pipeline}

Le calcul est découpé en cinq étapes indépendantes, chacune lisant la sortie de la précédente.
Chacune écrit un journal détaillé et dispose d'un auto-contrôle (\col{-{}-check}) qui rejoue sa
logique sur un jeu de données construit à la main, aux résultats connus d'avance.

\begin{center}
\small
\begin{tabular}{@{}l>{\RaggedRight}p{5.1cm}>{\RaggedRight}p{5.4cm}@{}}
\toprule
\textbf{Étape} & \textbf{Ce qu'elle fait} & \textbf{Ce qu'elle produit} \\
\midrule
01 & lit les quatre bases, calcule âge, scolarité, chômage, agrège au ménage &
\col{preconstruction\_...} (12\,965 $\times$ 38) \\
02 & applique le vecteur $z$ &
\col{vecteur\_z}, \col{matrice\_situationnelle\_...} (16 indicateurs 0/1) \\
03 & construit $w$, pondère, calcule $c_i$ &
\col{vecteur\_w}, \col{matrice\_privations\_ponderees} \\
04 & censure à $k$, calcule $c_i(k)$ &
\col{matrice\_privations\_censuree} \\
05 & calcule $H$, $A$, \ipm{}, leurs intervalles de confiance, les contributions et les
désagrégations (milieu, région, département, sous-préfecture, sexe et taille) --- pour les
\textbf{deux variantes}, nationale et internationale &
\col{indices\_ipm\_ci} et \col{indices\_ipm\_international}, leurs contributions et un
classeur \col{.xlsx} chacune \\
\bottomrule
\end{tabular}
\end{center}

\begin{verbatim}
python pipeline/orchestrateur.py            # les cinq étapes, environ 17 secondes
python pipeline/orchestrateur.py --check    # les auto-contrôles seuls
\end{verbatim}

Les poids changeant d'une variante à l'autre, l'étape 05 refait pour chacune la pondération et la
censure des étapes 03 et 04, au lieu de reprendre leur matrice --- qui ne vaut que pour la
variante nationale. Celle-ci est justement \emph{comparée} à la matrice de l'étape 04 par un
contrôle automatique : si les deux chemins de calcul divergeaient, l'étape s'arrêterait.

Une étape en échec interrompt la chaîne : mieux vaut aucune sortie qu'une sortie partielle. Toutes
les tables sont écrites simultanément au format Stata (\col{sorties/dta/}) et texte
(\col{sorties/csv/}) ; le classeur de restitution est dans \col{sorties/xlsx/}. Le contenu de ce
classeur, colonne par colonne, fait l'objet d'un document séparé.

\vfill
\noindent\rule{\linewidth}{0.4pt}\\[2pt]
\small
Documents associés : \emph{Dictionnaire des sorties} (chaque colonne du classeur) et
\emph{Note sur la dimension Santé} (le choix de l'insécurité alimentaire). Sources du pipeline
dans \col{PythonIPM/pipeline/}.

\end{document}
